\chapter{The Unruh Effect and Wedge Modular Theory}
\label{ch:curved-unruh-effect-modular-theory}

\ConventionsLorentzian

The Unruh effect is the thermal character of the Minkowski vacuum when
observables are restricted to a Rindler wedge and time evolution is generated
by Lorentz boosts.  The precise local-QFT theorem is the
Bisognano--Wichmann theorem; the accelerated-detector calculation is a
concrete consequence in free-field examples.

\section{Rindler Wedges and Boost Flow}

In \(\mathbb M^D\), define the right Rindler wedge
\[
  W_R=\{x\in\mathbb M^D:x^1>|x^0|\}.
\]
The boost preserving \(W_R\) is
\[
  x^0(\eta)=x^0\cosh\eta+x^1\sinh\eta,\qquad
  x^1(\eta)=x^1\cosh\eta+x^0\sinh\eta .
\]
Let \(K\) be the self-adjoint generator of this boost representation,
\[
  U(\eta)=\exp(\ii\eta K).
\]
For a uniformly accelerated observer with proper acceleration \(a\), the
Rindler time parameter is \(\eta=a\tau\), where \(\tau\) is proper time.

\section{KMS Statement}

Let \(\mathcal A(W_R)\) be the von Neumann algebra generated by observables
localized in \(W_R\).  The vacuum state restricted to \(\mathcal A(W_R)\) is
KMS at inverse temperature \(2\pi\) with respect to boost time if, for
\(A,B\in\mathcal A(W_R)\), the function
\[
  F_{A,B}(\eta)=\langle\Omega,A\,U(\eta)B\,U(\eta)^{-1}\Omega\rangle
\]
extends to the strip \(0<\operatorname{Im}\eta<2\pi\) and satisfies the KMS
boundary relation
\[
  F_{A,B}(\eta+2\pi\ii)
  =
  \langle\Omega,U(\eta)B\,U(\eta)^{-1}A\Omega\rangle .
\]
In proper time, the temperature is therefore
\[
  T_{\rm Unruh}=\frac{a}{2\pi}.
\]

\section{Quoted Bisognano--Wichmann Theorem}

\begin{quotedtheorem}[Bisognano--Wichmann theorem]
\label{thm:bisognano-wichmann-boundary}
For a Wightman or Haag--Kastler QFT satisfying the standard positivity,
covariance, locality, spectrum, vacuum cyclicity and separating hypotheses,
the modular automorphism group of \((\mathcal A(W_R),\Omega)\) is the Lorentz
boost flow preserving \(W_R\), with modular parameter normalized so that the
vacuum is KMS at inverse temperature \(2\pi\) for boosts.
\end{quotedtheorem}

This is the Bisognano--Wichmann theorem in its wedge-modular form.  It is not
a property of a detector model; it is a structural theorem about the pair
consisting of the wedge algebra and the vacuum vector.  Its proof uses
Tomita--Takesaki modular theory together with analyticity consequences of the
spectrum condition and wedge locality.  This chapter uses the result with its
hypotheses displayed and separates it from the explicit free-field detector
computation below.

\section{Free Scalar Check}

For a free massive scalar field in \(D=2\), the positive-frequency two-point
function is
\[
  W(x)=\int\frac{dp}{2\pi\,2E_p}\,e^{-\ii E_p x^0+\ii p x^1},
  \qquad E_p=\sqrt{p^2+m^2}.
\]
Along the uniformly accelerated trajectory
\[
  x^0(\tau)=a^{-1}\sinh(a\tau),\qquad
  x^1(\tau)=a^{-1}\cosh(a\tau),
\]
the invariant separation between two points is
\[
  -(x(\tau)-x(0))^2
  =
  \frac{4}{a^2}\sinh^2\!\left(\frac{a(\tau-\ii0)}{2}\right).
\]
The Wightman function along the trajectory is therefore a function of
\(\sinh^2(a(\tau-\ii0)/2)\).  Since
\[
  \sinh\!\left(\frac{a(\tau+\ii 2\pi/a)}{2}\right)
  =
  -\sinh\!\left(\frac{a\tau}{2}\right),
\]
the two-point function has imaginary proper-time periodicity
\(\beta=2\pi/a\) with the KMS ordering prescription.  This verifies the
thermal response for the free field.

\section{Detector Response}

An Unruh--DeWitt detector with energy gap \(E>0\) coupled linearly to the
field has transition rate proportional to
\[
  \dot P(E)
  =
  \int_{-\infty}^{\infty}d\tau\,
  e^{-\ii E\tau}
  W(x(\tau),x(0)).
\]
The KMS condition implies detailed balance
\[
  \dot P(-E)=e^{\beta E}\dot P(E),
  \qquad \beta=\frac{2\pi}{a}.
\]
Thus the detector response is thermal at temperature \(a/(2\pi)\) when the
coupling is weak, the detector is operated for a long time, and switching
effects have been separated from the stationary rate.

\begin{openproblem}[Interacting wedge thermality in examples]
\label{op:interacting-wedge-thermality-examples}
Develop explicit interacting examples where the wedge algebra, modular flow,
and detector response can be computed or bounded beyond formal perturbation
theory.  The target is to connect the abstract Bisognano--Wichmann theorem
with concrete nonperturbative correlation functions and operational detector
limits.
\end{openproblem}
